% TeX Macro for lecture notes
% dimensions
\abovedisplayskip=6pt plus 3pt minus 1pt
\belowdisplayskip=6pt plus 3pt minus 1pt
\abovedisplayshortskip=0pt plus 3pt
\belowdisplayshortskip=4pt plus 3pt
% A4 size paper
\hsize 159.2 true mm
\vsize 244 true mm
% \parindent 0pt
% \parskip is used only for cpnotes: 
\parskip 4pt plus 1pt minus 0.5pt
\overfullrule 0pt
% fonts
\font\titlefont=cmbx10 scaled 1440
\font\lecfont=cmbx12
\font\sc=cmcsc10
\font\ninerm=cmr9
\font\ninett=cmtt9
\def\h{\hskip 1cm}
\def\t{\hskip 6mm}
% Define the verbatim program listing macro \listing{filename}
%
\def\uncatcodespecials{\def\do##1{\catcode`##1=12 }\dospecials}
\def\listing#1{\par\begingroup\setupverbatim\input#1 \endgroup}
\def\setupverbatim{\tt 
  \def\par{\leavevmode\endgraf} \catcode`\`=\active
  \obeylines \uncatcodespecials \obeyspaces
  \everypar{\t}
}
{\obeyspaces\global\let =\ } % let active space = \control space
{\catcode`\`=\active \gdef`{\relax\lq}}
%
% verbatim text macro with smallskip before and after
% \ttbg
% verbatim text
% \tted
\def\ttbg{\par\smallskip\begingroup\parskip=0pt\setupverbatim\doverbatim}
{\catcode`\|=0 \catcode`\\=12 % | is temporary escape character
  |obeylines|gdef|doverbatim^^M#1\tted{#1|endgroup|smallskip}}
%
% one line verbatim
% \v`text`
\def\v{\begingroup\t\tt\uncatcodespecials
  \obeyspaces\doverbatt}
\newcount\balance
{\catcode`<=1 \catcode`>=2 \catcode`\{=12 \catcode`\}=12
  \gdef\doverbatt`<\balance=1\verbatimloop>
  \gdef\verbatimloop#1<\def\next<#1\verbatimloop>%
    \if#1`\advance\balance by-1
     \ifnum\balance=0\let\next=\endgroup\fi\fi\next>>

% the one line verbatim with comment in 9pt rm
%\l{\v`test`\m{comment}}
\def\m#1{\hfill {\ninerm #1} \t}
\def\l{\hbox to \hsize}
% one line of tt font without comment
\def\ll#1{\leftline{\t\tt #1}}
\def\ttsp{\hphantom{2}}
\newskip\hugeskipamount
\hugeskipamount = 20mm plus3mm minus3mm
\def\hugeskip{\vskip\hugeskipamount}
\def\hugebreak{\par \ifdim\lastskip<\hugeskipamount
  \removelastskip \penalty-500 \hugeskip \fi}
\def\lecture#1{\hugebreak {\noindent\lecfont Lecture #1} \bigskip\nobreak}
\def\section#1{\medbreak {\noindent\sc #1} \smallskip\nobreak}
% macro for a book-like lecture notes
\newcount\chapno
\newcount\secno
\newcount\subsecno
\newcount\equationno
\global\chapno=0
\global\equationno=0
\outer\def\chap#1{\global\advance\chapno by 1 \global\secno=0
  \centerline{\lecfont Chapter \the\chapno\ \ #1}
  \par\bigskip\nobreak}
\outer\def\sec#1{\medbreak \global\advance\secno by 1 \global\subsecno=0
  \medskip\vbox{\noindent{\sc \the\secno.\ #1}\medskip}
  \nobreak}
\outer\def\subsec#1{\smallbreak \global\advance\subsecno by 1
  \medskip\vbox{\noindent{\sc \the\secno--\the\subsecno.\ #1}
  \smallskip}\nobreak}
\def\EQNO{\global\advance\equationno by 1 \eqno(\the\equationno) }

\def\sms{\smallskip}
\def\fninety{$^{\scriptscriptstyle 90}$}

%def \problem and \answer with automatic numbering
\newcount\probno
\probno=0

\outer\def\problem{\medbreak\smallskip \global\advance\probno by 1
\item{{\bf\the\probno}.}}

\chardef\other=12
\newwrite\ans
\immediate\openout\ans=answers %file for answers to problems
\outer\def\answer{\par\medbreak
  \immediate\write\ans{}
  \immediate\write\ans{\string\ansno{\the\probno}}
  \copytoblankline}
\def\copytoblankline{\begingroup\setupcopy\copyans}
\def\setupcopy{\def\do##1{\catcode`##1=\other}\dospecials
    \catcode`\|=\other \obeylines}
{\obeylines \gdef\copyans#1
  {\def\next{#1}%
  \ifx\next\empty\let\next=\endgroup %
  \else\immediate\write\ans{\next} \let\next=\copyans\fi\next}}
\def\ansno#1{\medbreak\noindent{\bf #1:\ }}

\input epsf.tex
%end of macro definitions

